\section{Results}
\label{sec:results}

All results in this section refer to the final chronological test split. Model
rankings are stated with the metric used for the ranking. Positive MAE gain
means that the feature-based model has lower MAE than the last-hour baseline;
positive ablation $\Delta$MAE means that removing the feature or feature family
increases MAE.

The pedestrian-flow task is the only target family for which feature-based
models lower one-hour MAE and WAPE relative to immediate persistence for both
targets. Table~\ref{tab:flow-results} reports the tree and boosting alternatives
with the lowest one-hour MAE, together with last-hour persistence.

\begin{table}
\caption{One-hour pedestrian-flow model leaderboard ranked by MAE on the final chronological test split.}
\label{tab:flow-results}
\centering
\scriptsize
\setlength{\tabcolsep}{3.5pt}
\begin{tabular}{llrrrrr}
\toprule
Target & Model & MAE & $R^2$ & WAPE & Fit s & ms/row \\
\midrule
ingressi\_borgo & Extra Trees & 18.69 & 0.554 & 37.81 & 0.50 & 0.119 \\
ingressi\_borgo & Random Forest & 19.23 & 0.545 & 38.90 & 0.68 & 0.080 \\
ingressi\_borgo & XGBoost & 19.89 & 0.539 & 40.24 & 0.31 & 0.005 \\
ingressi\_borgo & LightGBM & 20.98 & 0.524 & 42.42 & 0.17 & 0.006 \\
ingressi\_borgo & HGB & 21.16 & 0.508 & 42.80 & 0.43 & 0.014 \\
ingressi\_borgo & Last-hour & 21.29 & 0.636 & 43.06 & 0.00 & 0.010 \\
uscite\_borgo & Extra Trees & 15.39 & 0.826 & 28.93 & 0.50 & 0.099 \\
uscite\_borgo & Random Forest & 15.54 & 0.828 & 29.20 & 0.61 & 0.100 \\
uscite\_borgo & XGBoost & 15.66 & 0.821 & 29.42 & 0.31 & 0.005 \\
uscite\_borgo & HGB & 16.07 & 0.823 & 30.19 & 0.33 & 0.014 \\
uscite\_borgo & LightGBM & 16.22 & 0.807 & 30.49 & 0.17 & 0.006 \\
uscite\_borgo & Last-hour & 20.54 & 0.719 & 38.61 & 0.00 & 0.006 \\
\bottomrule
\end{tabular}
\end{table}

For entrances, Extra Trees reduces MAE by 12.2\% and WAPE by 5.3 percentage
points relative to last-hour persistence. For exits, it reduces MAE by 25.1\%
and WAPE by 9.7 percentage points. The $R^2$ ordering is different: last-hour
persistence has the highest $R^2$ for entrances, and Random Forest has the
highest $R^2$ for exits. The supported claim is therefore lower absolute error
for the Extra Trees flow forecasts, not universal dominance on every metric.
The close MAE values for Extra Trees, Random Forest, and XGBoost show that the
one-hour flow result is not tied to a single tree implementation.

The selected one-hour feature sets also differ by target family
(Table~\ref{tab:feature-groups}). Flow uses mostly recent camera-history
features, plus weather, calendar, TIM, and one target-history feature.
Nationality uses TIM, event, weather, and target-history features. Age uses TIM,
target-history, and one weather feature. Event variables are not selected for
one-hour flow or one-hour age models.

\begin{table}
\caption{One-hour selected-feature group composition, reported as counts of selected predictors by feature family.}
\label{tab:feature-groups}
\centering
\scriptsize
\begin{tabular}{lrrrrrrr}
\toprule
Track & Calendar & Events & Weather & Pedestrian & Target hist. & TIM & Total \\
\midrule
Flow & 2 & 0 & 6 & 14 & 1 & 2 & 25 \\
Nationality & 0 & 3 & 2 & 0 & 3 & 7 & 15 \\
Age & 0 & 0 & 1 & 0 & 8 & 6 & 15 \\
\bottomrule
\end{tabular}
\end{table}

The multi-horizon comparison in Table~\ref{tab:flow-horizon} limits the flow
claim to short horizons. Feature-based models improve over last-hour
persistence for both flow targets at one hour. At three hours only entrances
show a positive MAE gain. At six hours entrances improve by 8.4\%, while exits
are effectively tied with persistence (+0.2\% MAE gain). At 12 and 24 hours,
last-hour persistence has lower MAE than the best feature-based model.

\begin{table}
\caption{Pedestrian-flow comparison between last-hour persistence and the lowest-MAE feature-based model across horizons. Positive gain means lower MAE than persistence.}
\label{tab:flow-horizon}
\centering
\scriptsize
\begin{tabularx}{\textwidth}{rlXrrrr}
\toprule
Horizon & Target & Feature model & Last-hour MAE & Feature MAE & Feature $R^2$ & MAE gain \\
\midrule
1 h & ingressi\_borgo & Extra Trees & 21.29 & 18.69 & 0.554 & +12.20\% \\
1 h & uscite\_borgo & Extra Trees & 20.54 & 15.39 & 0.826 & +25.07\% \\
3 h & ingressi\_borgo & Random Forest & 21.39 & 21.13 & 0.438 & +1.22\% \\
3 h & uscite\_borgo & Random Forest & 20.63 & 20.96 & 0.611 & -1.62\% \\
6 h & ingressi\_borgo & Extra Trees & 21.56 & 19.76 & 0.558 & +8.37\% \\
6 h & uscite\_borgo & Extra Trees & 20.78 & 20.73 & 0.611 & +0.21\% \\
12 h & ingressi\_borgo & Extra Trees & 21.44 & 24.15 & 0.441 & -12.64\% \\
12 h & uscite\_borgo & Extra Trees & 20.75 & 26.87 & 0.443 & -29.51\% \\
24 h & ingressi\_borgo & XGBoost & 21.73 & 22.59 & 0.471 & -3.98\% \\
24 h & uscite\_borgo & XGBoost & 20.84 & 23.81 & 0.526 & -14.23\% \\
\bottomrule
\end{tabularx}
\end{table}

Figure~\ref{fig:flow-horizon-gain} visualizes the same comparison as a signed
MAE reduction relative to last-hour persistence. The figure is retained because
it makes the non-monotonic horizon behavior visible: the feature-based advantage
is largest at one hour, partly reappears for entrances at six hours, and is
negative at longer horizons.

\begin{figure}
\centering
\includegraphics[width=0.78\textwidth]{feature_gain_vs_last_hour_by_horizon_flow.pdf}
\caption{Pedestrian-flow MAE reduction of the lowest-MAE feature model relative to last-hour persistence. Positive values mean lower MAE than persistence.}
\label{fig:flow-horizon-gain}
\end{figure}

Rolling validation gives a stricter temporal robustness check than the final
split alone. Table~\ref{tab:rolling-validation} reports the lowest-MAE model
averaged over three rolling folds at the one-hour horizon. For exits, Extra
Trees is also the rolling-validation winner. For entrances, rolling validation
selects last-hour persistence even though Extra Trees has lower MAE on the
final September test split. The entrance result is therefore a final-period
error reduction, while the exit result is supported by both rolling validation
and the final split.

\begin{table}
\caption{One-hour rolling-validation winners by target. Values are means over three chronological folds.}
\label{tab:rolling-validation}
\centering
\scriptsize
\setlength{\tabcolsep}{3.5pt}
\begin{tabular}{lllrrrr}
\toprule
Track & Target & Rolling winner & Folds & MAE mean & MAE std & $R^2$ mean \\
\midrule
Flow & entrances & Last-hour & 3 & 17.35 & 2.81 & 0.631 \\
Flow & exits & Extra Trees & 3 & 15.95 & 3.31 & 0.713 \\
Nationality & Italian & Last-hour & 3 & 102.38 & 26.85 & 0.886 \\
Nationality & foreign & Last-hour & 3 & 10.34 & 3.40 & 0.866 \\
Age & $<18$ & Last-hour & 3 & 5.15 & 0.97 & 0.711 \\
Age & 18--30 & Last-hour & 3 & 18.22 & 5.50 & 0.802 \\
Age & 31--40 & Last-hour & 3 & 18.49 & 4.41 & 0.833 \\
Age & 41--50 & Last-hour & 3 & 24.74 & 6.75 & 0.819 \\
Age & 51--60 & Last-hour & 3 & 29.07 & 5.60 & 0.841 \\
Age & 61+ & Last-hour & 3 & 28.64 & 4.57 & 0.809 \\
\bottomrule
\end{tabular}
\end{table}

TIM nationality and age targets have a different behavior. Table~\ref{tab:persistence}
shows that, at one hour, last-hour persistence has lower MAE than the
lowest-MAE feature-based model for all eight TIM targets. The same pattern holds across
all 40 TIM target-horizon cases in the complete multi-horizon target summary.
Feature-rich models often achieve
positive $R^2$, but their absolute error is higher than immediate persistence.
For these TIM indicators, the predictable component in this experiment is
therefore dominated by autoregressive continuity.

\begin{table}
\caption{Persistence versus the lowest-MAE feature-based model at the one-hour horizon for TIM targets.}
\label{tab:persistence}
\centering
\scriptsize
\begin{tabularx}{\textwidth}{llrrXrr}
\toprule
Track & Target & Last-hour MAE & Last-hour $R^2$ & Best feature model & Feature MAE & Feature $R^2$ \\
\midrule
Nationality & tim\_Ni\_mean15 & 97.80 & 0.858 & ridge & 157.88 & 0.604 \\
Nationality & tim\_Ns\_mean15 & 8.00 & 0.815 & Extra Trees & 12.00 & 0.626 \\
Age & tim\_F1\_mean15 & 4.91 & 0.711 & two-stage ridge & 10.13 & 0.191 \\
Age & tim\_F2\_mean15 & 15.26 & 0.787 & ridge & 21.14 & 0.457 \\
Age & tim\_F3\_mean15 & 16.60 & 0.800 & ridge & 25.02 & 0.488 \\
Age & tim\_F4\_mean15 & 21.88 & 0.792 & two-stage ridge & 34.45 & 0.446 \\
Age & tim\_F5\_mean15 & 27.47 & 0.848 & Extra Trees & 46.15 & 0.590 \\
Age & tim\_F6\_mean15 & 26.23 & 0.732 & HGB & 37.20 & 0.506 \\
\bottomrule
\end{tabularx}
\end{table}

Group-level ablation gives a coarse interpretation of the best feature-based
models. Table~\ref{tab:ablation} reports the feature family whose removal gives
the largest average one-hour MAE increase within each track. For flow, the
dominant family is pedestrian history. For nationality and age, the dominant
families are TIM-derived variables and target history. Weather and event
features can be selected, but their one-hour group-level effect is smaller than
the dominant history or TIM groups.

\begin{table}
\caption{Largest average MAE degradation under one-hour leave-one-group-out ablation. Positive values mean worse performance after removing the feature family.}
\label{tab:ablation}
\centering
\begin{tabular}{lcr}
\toprule
Track & Highest-impact removed group & Mean $\Delta$MAE \\
\midrule
Flow & pedestrian history & +6.01 \\
Nationality & TIM & +67.41 \\
Age & TIM & +9.20 \\
\bottomrule
\end{tabular}
\end{table}

The leave-one-feature-out ablation in Table~\ref{tab:single-feature-ablation}
is narrower than group ablation: it removes one selected feature at a time and
re-trains the lowest-MAE feature-based model for the affected target, without
retuning. It identifies local dependencies inside the selected feature set. For
one-hour flow, the largest single-feature effects are small in count units
relative to the group-level pedestrian-history effect. For TIM targets, the
largest single-feature effects are usually TIM or target-history lags, which is
consistent with the persistence-dominated performance comparison.

\begin{table}
\caption{Largest leave-one-feature-out MAE degradation by track and horizon, after re-training the selected feature-based model without one predictor.}
\label{tab:single-feature-ablation}
\centering
\scriptsize
\setlength{\tabcolsep}{2.5pt}
\begin{tabularx}{\textwidth}{rl l l X l r}
\toprule
Horizon & Track & Target & Model & Removed feature & Group & $\Delta$MAE \\
\midrule
1 h & Flow & \path|ingressi_borgo| & Extra Trees & \path|hour| & calendar & +0.85 \\
1 h & Nationality & \path|tim_Ni_mean15| & Ridge & \path|tim_Tc_mean15_lag1h| & TIM & +38.56 \\
1 h & Age & \path|tim_F6_mean15| & HGB & \path|tim_F6_mean15_lag1h| & target hist. & +12.96 \\
3 h & Flow & \path|uscite_borgo| & Random Forest & \path|entra_piazzarocca_ingresso_lag1h| & pedestrian & +0.65 \\
3 h & Nationality & \path|tim_Ni_mean15| & Extra Trees & \path|tim_Ns_mean15_lag1h| & target hist. & +19.29 \\
3 h & Age & \path|tim_F6_mean15| & Poisson & \path|tim_F6_sum15_roll24h_mean| & target hist. & +20.48 \\
6 h & Flow & \path|uscite_borgo| & Extra Trees & \path|tim_Vp_max15_roll6h_mean| & TIM & +1.88 \\
6 h & Nationality & \path|tim_Ni_mean15| & Ridge & \path|tim_F6_mean15_lag24h| & TIM & +21.53 \\
6 h & Age & \path|tim_F6_mean15| & Poisson & \path|tim_F6_sum15_roll24h_mean| & target hist. & +49.14 \\
12 h & Flow & \path|uscite_borgo| & Extra Trees & \path|tim_Vi_mean15_lag24h| & TIM & +0.95 \\
12 h & Nationality & \path|tim_Ni_mean15| & Poisson & \path|tim_Gf_mean15_roll24h_mean| & TIM & +238.03 \\
12 h & Age & \path|tim_F6_mean15| & HGB & \path|tim_F6_sum15_roll24h_mean| & target hist. & +104.55 \\
24 h & Flow & \path|uscite_borgo| & XGBoost & \path|tim_Vp_mean15_lag24h| & TIM & +4.95 \\
24 h & Nationality & \path|tim_Ni_mean15| & XGBoost & \path|tim_Gf_mean15_lag24h| & TIM & +16.66 \\
24 h & Age & \path|tim_F6_mean15| & XGBoost & \path|tim_F6_mean15_lag24h| & target hist. & +13.28 \\
\bottomrule
\end{tabularx}
\end{table}

The event-feature analysis is mixed. The event-aware dataset contains 42 event
features, and at least one event signal is non-zero in 2,132 of 2,735 modeling
hours. Table~\ref{tab:events} shows that event variables enter the selected
feature set mainly for nationality and age at 3--12 hours; they enter the flow
feature set only at six hours. The leave-one-group-out event effect is positive
for nationality at one, three, and six hours and for flow at six hours, but it
is negative for nationality at 12 and 24 hours and for age at 3, 6, and 12
hours. This is evidence of predictive association in some fitted models, not a
stable performance gain.

\begin{table}
\caption{Event-feature selection and leave-one-group-out event ablation by horizon and track. Positive $\Delta$MAE means removing event features worsens the selected feature-based model.}
\label{tab:events}
\centering
\scriptsize
\begin{tabular}{rlrrr}
\toprule
Horizon & Track & Event features & Event share & Mean event $\Delta$MAE \\
\midrule
1 h & flow & 0/25 & 0.0\% & -- \\
1 h & nationality & 3/15 & 20.0\% & +3.39 \\
1 h & age & 0/15 & 0.0\% & -- \\
3 h & flow & 0/15 & 0.0\% & -- \\
3 h & nationality & 3/15 & 20.0\% & +2.53 \\
3 h & age & 4/15 & 26.7\% & -0.04 \\
6 h & flow & 3/15 & 20.0\% & +0.64 \\
6 h & nationality & 5/15 & 33.3\% & +0.60 \\
6 h & age & 7/15 & 46.7\% & -0.92 \\
12 h & flow & 0/15 & 0.0\% & -- \\
12 h & nationality & 6/15 & 40.0\% & -20.13 \\
12 h & age & 8/15 & 53.3\% & -4.61 \\
24 h & flow & 0/15 & 0.0\% & -- \\
24 h & nationality & 2/15 & 13.3\% & -0.20 \\
24 h & age & 0/15 & 0.0\% & -- \\
\bottomrule
\end{tabular}
\end{table}

SHAP summaries support the same cautious interpretation. Some event variables
rank highly for individual targets: sport/motor event intensity is the top SHAP
feature for \path|tim_Ni_mean15| at three hours, for \path|tim_F5_mean15| and
\path|tim_F6_mean15| at six hours, and for \path|tim_F2_mean15| and
\path|tim_F5_mean15| at 12 hours. Those target-level rankings do not overturn
the ablation result: the event signal is visible in parts of the demographic
models, but the current event encoding does not produce a consistent reduction
in test error across tracks and horizons.

The main operational result is therefore narrow and explicit. For short-horizon
pedestrian-flow forecasting, recent camera dynamics plus selected contextual
covariates provide lower MAE and WAPE than immediate persistence at one hour.
For TIM nationality and age indicators, immediate persistence remains the
lowest-MAE forecasting baseline in this dataset. Event variables should be kept
as exploratory covariates and diagnostics, not presented as a validated driver
of forecast accuracy.
